\cotefile{157-3}
\scanpage{252}
\mspage{202}

\[
(\unsure{\xi|\zeta}*\xi')*\xi''\ \unsure{\xrightarrow{\sim}|\overset{\sim}{\text{---}}}\ \xi*(\xi'*\xi'')
\]
Il \deleted{\illegible{1}} y a deux façons d'\unsure{immerger|envoyer} un produit binaire
\[
\xi*\xi',\quad \xi'*\xi'',\quad \xi*\xi''
\]
dans le produit ternaire, \unsure{eh|et} bien elles doivent \deleted{\unsure{dans}} chacun des trois cas, être égales. Et de \unsure{$=$|même} il y a, ceci admis, \deleted{\unsure{deux une}} deux façons d'\unsure{immerger|envoyer} $\xi$ dans le produit ternaire, en factorisant soit par $\xi*\xi'$, soit par $\xi*\xi''$, et de $=$ pour les deux autres facteurs $\xi',\xi''$; et là encore, ces deux façons doivent être égales, dans chacun des \unsure{ces|en} trois cas -- donc $3+3=6$ égalités en \unsure{tout} à écrire \inserted{et à vérifier} -- on a vu pire !

La philosophie \unsure{calculatoire} devrait être celle-ci~: \unsure{Si on a|Lié à} un produit $n$-uple
\[
\xi_1*\xi_2*\text{---}*\xi_n
\]
\unsure{(ici|(où} je me \unsure{dispense} de mettre des \unsure{parenthèses} ....), pour \unsure{tout} produit $m$-uple partiel ($m\in[1,n-1]$)
% archive sheet number: 252
\scanpage{253}
\mspage{203}

\[
\unsure{\xi|\zeta}_{i_1}*\xi_{i_2}*\cdots*\xi_{i_m}\qquad(i_1<i_2\ \unsure{\ldots|{-}}\ <i_m),
\]
en utilisant à tour de bras les \unsure{isom.|isos.} d'associativité et les $\alpha_{\xi,\xi'},\ \beta_{\xi,\xi'}$, il n'y a \unsure{pourtant|peut-être} qu'une seule façon de l'\unsure{envoyer|immerger} dans le produit $n$-uple. J'ai bien l'impression que ça doit être vrai.

(Est-ce que ça implique déjà que $\unsure{\mathcal{C}|C}$ est $1$-connexe ?)

En fait, il y a bien une unité bilatère, mais elle n'est pas dans $\mathcal{C}$, mais dans la catégorie plus \unsure{grande} \deleted{\unsure{de}}
\[
{}_{\mathrm{id}_{\mathrm{Cat}_{/S}}}\!\backslash\,{}^{\underline{\mathrm{Hom}}(\mathrm{Cat}_{/S},\mathrm{Cat}_{/S})}\ ,
\]
qui elle, comme \deleted{\unsure{toute}} catégorie
\[
{}_{\mathrm{id}_{\unsure{\mathcal{X}|\mathscr{X}}}}\!\backslash\,{}^{\underline{\mathrm{Hom}}(\mathcal{X},\mathcal{X})}\ ,
\]
est une catégorie monoïdale, i.e. un produit associatif \uline{et} unitaire. Les \deleted{\unsure{compatib.}} flèches $\alpha_{\xi,\xi'},\ \beta_{\xi,\xi'}$ sont les flèches qui résultent de la présence de cette unité. (vu que $\mathcal{C}$ est une sous-catégorie pleine), et \unsure{donc|alors} il doit devenir clair qu'on \unsure{a} les compatibilités dites
% archive sheet number: 253
\scanpage{254}
\mspage{204}

. Mais même une catégorie monoïdale, est elle $1$-connexe ? Par exemple $B_G$, pour $G$ un groupe \unsure{disons|discret}, est bien une catégorie monoïdale et même ``groupale'', une \unsure{Gr}-catégorie (tout ce qu'il y a de strict), et pourtant elle n'est \uline{pas} $1$-connexe !

Bien sûr, on peut court-circuiter la question de la $1$-connexité de $\unsure{\mathcal{C}|C}$, en notant qu'après \deleted{tout} \unsure{tout} on a déjà a priori un objet privilégié dans
\[
\underline{\mathrm{Hom}}(\mathrm{Cat}_{/S},\mathrm{HOT}_{W}^{\unsure{\mathrm{lc}|\mathrm{le}}}(S))
\]
\unsure{i.e.|scil.} un foncteur canonique
\[
\mathrm{Cat}_{/S}\longrightarrow\mathrm{HOT}_{W}^{\mathrm{lc}}(S),
\]
savoir \deleted{\unsure{celui}} le foncteur \deleted{\unsure{quasi-inverse}}\inserted{composé} \deleted{\unsure{du} foncteur}
\[
\deleted{\mathrm{HOT}^{\mathrm{lc}}_W(S)\overset{\mathrm{def}}{=}(\underline{\mathrm{Fib}}_W S)(W_S^{\unsure{f|t}})^{-1}\longrightarrow\mathrm{Cat}}
\]
\diagram{\begin{tikzcd}[ampersand replacement=\&,column sep=small,row sep=small]
\text{\unsure{(x)|(*)}}\quad \mathrm{Cat}_{/S}
  \arrow[r,-{Bar},"{\substack{\text{localisation}\\\text{par }W_S}}"']
  \&(\mathrm{Cat}_{/S})W_S^{-1}\arrow[r]\&\mathrm{HOT}^{\mathrm{lc}}_W(S)\arrow[d,equal,"{\mathrm{def}}"]\\
\&\&(\underline{\mathrm{Fib}}_W S)(W_S^{\unsure{t|f}})^{-1}
\end{tikzcd}}
où la deuxième flèche est définie comme le quasi-inverse de
\[
\text{\unsure{(xx)|(**)}}\quad(\underline{\mathrm{Fib}}_W S)(W_S^{t})^{-1}\longrightarrow(\mathrm{Cat}_{/S})W_S^{-1}
\]
% archive sheet number: 254
\scanpage{255}
\mspage{205}

déduit de l'inclusion
\[
(\underline{\mathrm{Fib}}_W S,W_S^{\unsure{t|f}})\lhook\joinrel\longrightarrow(\mathrm{Cat}_{/S},W_S).
\]
Le fait que \unsure{(xx)|ceci} soit une équivalence, se prouve justement par le fait de l'existence d'un foncteur $C:\mathrm{Cat}_{/S}\to\mathrm{Cat}_{/S}$ et d'un $i:\mathrm{id}_{\mathrm{Cat}_{/S}}\to C$, ayant les propriétés que \unsure{l'on|on} sait, i.e. du fait que $\unsure{\mathcal{C}|C}\ne\varnothing$. De plus, la donnée d'un objet $C$ de $\mathcal{C}$ donne un moyen commode d'expliciter le quasi-inverse du dit foncteur. Mais ce quasi-inverse \unsure{à|a} une existence canonique \unsure{etc|et} indépendante du choix d'un $C$, i.e. d'un mode de calcul. Le fait que le foncteur composé \inserted{$\unsure{\varphi|\phi}$} dans \unsure{(x)|(*)} puisse se ``relever'' \inserted{(à isom près)} \unsure{en|via} un foncteur
\[
C:\mathrm{Cat}_{/S}\longrightarrow\underline{\mathrm{Fib}}_W S,
\]
et que l'isomorphisme \inserted{entre $\varphi(X)\unsure{\flat|b|!}$ et $\mathrm{hot}_{W,S}(C(X))$} \margin{$\flat=\mathrm{hot}_{W,S}(X)$} \unsure{puisse s'expliciter|en question s'explicite}, lui \unsure{même|aussi}, par une flèche $X$-fonctorielle
\[
i_X:X\longrightarrow C(X)\qquad i_X\in W
\]
dans $\mathrm{Cat}_{/S}$ lui-\unsure{$=$|même} (et non \inserted{seulement} dans le \unsure{localisé} $(\mathrm{Cat}_{/S},W_S^{-1})$,
% archive sheet number: 255
\scanpage{256}
\mspage{206}

est une vertu particulière des \unsure{Vogons|Vosges|Vagos}, qui ne doit pas pour autant faire oublier le caractère intrinsèque des flèches \unsure{(x)|(*)}, indépendamment de \unsure{tout} mode de calcul.

\margin{\rule{0.35pt}{4\baselineskip}\hspace{2pt}\rule{0.35pt}{4\baselineskip}}%
La clef de la situation, et des questions de canonicité, me semble être dans le fait que \unsure{(xx)|(**)} est une équivalence de catégories -- et il faut avouer que, tout familier et quasiment banal qu'il soit finalement devenu pour moi, c'est là un fait tout \unsure{à|:} fait remarquable et nullement évident a priori. (Mais sans doute peut-on \unsure{le avoir|l'avoir} aussi par les \unsure{fourbis} généraux à la Quillen-Thomason, sans se battre avec un formalisme des chemins ...) Les ``isom can'' qu'on peut établir entre une ``théorie $C$'' et une ``théorie $C'$'' (via deux constructions différentes, p. ex, de catégories de chemins, ... ou via l'approche de Quillen-Thomason, sans chemins) doivent toujours revenir à \unsure{celui|cela|celui-ci}
% archive sheet number: 256
\scanpage{257}
\mspage{207}

qu'on peut définir via la comparaison de la ``théorie $C$'' d'une part, la ``théorie $C'$'' de l'autre, avec ce que nous donne directement la flèche composée \unsure{(x)|(*)}. (p. 204).

Finalement, il me semble que l'\uline{existence} d'une théorie des ``fibres homotopiques'' (\unsure{i.e.|. . .} ``$W$-homotopiques'') d'une flèche $X\to S$ dans $\mathrm{Cat}$, résulte de l'équivalence \unsure{(xx)|(**)} (p. 204), et \inserted{cette théorie} pourrait sans doute dans une certaine mesure se développer à partir de là, sans choix d'un foncteur $C$ ni d'une structure de ``catégorie de modèles au sens de Quillen\unsure{''|,}

\unsure{9.|4.} \uline{Axiomatisation des ``fibres \inserted{$W$-}homotopiques.''} La réflexion de hier me fait revenir sur l'axiomatique dans des ``cat. de modèles''
\[
\unsure{(\,.\quad {}'|(\,.\quad .}\qquad(M,W)
\]
% slanted
\margin{$W$ satisfait l'\unsure{``|{}}axiome de saturation faible~: les iso dans $W$, et si deux parmi $f,g,gf$ dans $W$, le troisième \unsure{itou|aussi} \unsure{(si)|(sic)|bien sûr)}.}
plus générales que $\mathrm{Cat}$, du $1^{\mathrm{er}}$ \unsure{append.} de XII (\unsure{cf th.-scholie 3|cf th.-scholie 2} p. 108--113) Pour le moment, on va supposer seulement $M$ stable par produits fibrés (on n'a pas besoin d'objet final). On introduit l'ensemble \deleted{\unsure{des} $\mathrm{Fib}_W\subset\mathrm{Fl}(M)$}
\[
\mathrm{Fib}_W\subset\mathrm{Fl}(M),
\]
% archive sheet number: 257
\scanpage{258}
\mspage{208}

des \uline{$W$-fibrations} (sous-entendu ``\unsure{ultra-faibles|ultra. faibles}'' !) par la condition

$f:X\to Y$ est dans $\mathrm{Fib}_W$ ssi \deleted{pour tout} le foncteur chgt de base par $X$, $Y'\mapsto X'=X\times_Y Y'$; \deleted{\unsure{de $M$}},
\[
M_{/Y}\longrightarrow M_{/X}
\]
transforme $\deleted{\unsure{Y}}W_X$ dans $W_Y$,

\deleted{\unsure{$W$.}} \quad $W_X\subset\mathrm{Fl}(M_{/X})$ défini par
\[
u\in W_X\Longleftrightarrow\sigma(u)\in W
\]
($\sigma:M_{/X}\to M$ le foncteur ``source'', ou ``oubli de $X$'')
% slanted
\margin{\uline{NB} Donner ici déjà la prop.~: \uline{$\mathrm{Fib}_W$ stable par composition et par chgt de base}, \uline{et contient les iso}, stable par \unsure{fact. dir.}}

Posons, pour tout $S$ dans $M$
\[
\mathrm{Ho}^{!}_W(S)=(M_{/S})W_S^{-1},
\]
\deleted{\unsure{Qui|qui}} joue le rôle d'une catégorie de $W$-types d'homotopie relatifs \uline{loc\textsuperscript{t} constants} sur $S$.

On aimerait décrire une loi contravariante \deleted{\unsure{pour}} $S$ pour $\mathrm{Ho}^{!}_W(S)$, i.e. pour
\[
f:S'\longrightarrow S
\]
flèche quelconque dans $M$, définir
\[
f^{*}_W:\mathrm{Ho}^{!}_W(S)\unsure{\xrightarrow{?}|\overset{?}{\text{---}}}\mathrm{Ho}^{!}_W(S').\tag*{\unsure{(x)|(*)}}
\]
La difficulté est, bien sûr, que le foncteur changement de base
\[
f^{*}:M_{/S}\longrightarrow M_{/S},
\]
n'est \uline{pas} en général compatible avec les
% archive sheet number: 258
\scanpage{259}
\mspage{209}

localiseurs $W_S,W_{S'}$. La \deleted{\unsure{façon}} notion de $W$-fibration sert donc en premier lieu à dégager la condition la plus évidente sur une flèche $f:S'\to S$ dans $M$, pour que \deleted{\unsure{elle}} le foncteur \unsure{changement|chgt} de base \deleted{\unsure{induit}} soit compatible avec les localiseurs $W_S,W_{S'}$ (c'est \uline{ça} la définition), donc induise $(*)$ sur les \unsure{localisées}.

Le ``\uline{problème des \deleted{\unsure{types}} fibres homotopiques}'' est justement de définir \unsure{(x)|(*)} pour \uline{toute} \deleted{\unsure{classe}}\inserted{flèche} dans $M$, pas seulement pour les $W$-fibrations, de telle façon que les $\mathrm{Ho}^{!}_W(S)$ forment, pour $S$ variable, les fibres d'une catégorie fibrée remarquable sur $M$. Quand $M$ admet un objet final \inserted{$e$}, \unsure{de} sorte que l'on peut regarder les flèches
\[
p:e\unsure{\longrightarrow|\text{---}}S\qquad\text{($p$ comme ``point'')}
\]
dans $M$ comme des ``points'' de $S$, la ``fibre $W$-homotopique'' d'un objet quelconque $X/S$ de $\deleted{\unsure{\mathrm{Cat}|C}}\unsure{M}_{/S}$ en ledit point sera par définition l'élément de
\[
\mathrm{Ho}_W=\mathrm{Ho}_W(M)\simeq\mathrm{Ho}^{!}_W(e)
\]
défini par
\[
p^*(\mathrm{ho}^{!}_{W/S}(X/S)),
\]
% archive sheet number: 259
\scanpage{260}
\mspage{210}

où
\[
\mathrm{ho}^{!}_{W/S}:M_{/S}\longrightarrow\mathrm{Ho}^{!}_W(S)\overset{\mathrm{def}}{=}(M_{/S})(W_S)^{-1}
\]
est le foncteur de localisation canonique.

Pour \deleted{tout} $S$ dans $M$, soit
\[
\underline{\mathrm{Fib}}_W S\subset M_{/S}
\]
la sous-catégorie pleine formée des objets $X_{/S}$ de $M_{/S}$ tels que $X\to S$ soit $W$-fibrant.

On va définir une variante
\[
\mathrm{Ho}^{\unsure{\mathrm{lc}|\mathrm{le}}}_W(S)
\]
\deleted{\unsure{comme|covariante}} de $\mathrm{Ho}^{!}_W(S)$, comme
\[
\mathrm{Ho}^{\mathrm{lc}}_W(S)\overset{\mathrm{def}}{=}(\underline{\mathrm{Fib}}_W S)(W_S^{u})^{-1}
\]
où $W_S^{u}\subset\deleted{\unsure{\mathrm{Fib}}}\mathrm{Fl}(\underline{\mathrm{Fib}}_W S)$ est le localiseur formé des flèches \inserted{$u:X\to Y$,} qui sont ``universellement \inserted{sur $S$} dans $W$'', i.e telles que pour toute flèche $S'\to S$ de $\deleted{\unsure{\mathrm{Cat}|C}}\unsure{M}$ de but $S$, la flèche déduite par chgt. de base,
\[
u':X'=X\times_S S'\longrightarrow Y'=Y\times_S S'
\]
est dans $W_{S'}$.

On a ``fait ce qu'il fallait'' pour que $\mathrm{Ho}^{\mathrm{lc}}_W(S)$ soit contravariant en $S$ pour des flèches \uline{quelconques} \inserted{$f:S'\to S$} de $M$, elles définissent
\[
f_{\deleted{\unsure{*}}W}^{*\,\mathrm{lc}}:\mathrm{Ho}^{\mathrm{lc}}_W(S)\longrightarrow \mathrm{H}\deleted{\unsure{\mathrm{O}}}\mathrm{o}^{\mathrm{lc}}_W(S'),\tag*{\unsure{(xx)|(**)}}
\]
% archive sheet number: 260
\scanpage{261}
\mspage{\deleted{\unsure{2|1}}211}

puisque les $W$-fibrations sont stables par changement de base, et que par suite $f^*: \deleted{\unsure{\mathrm{Cat}|C}}\unsure{M}_{/S}\unsure{\text{---}|\to}M_{/S'}$ induit
\[
f^*_{\underline{\mathrm{Fib}}_W}:\underline{\mathrm{Fib}}_W S\ \deleted{\unsure{M}}\unsure{\text{---}|\longrightarrow}\underline{\mathrm{Fib}}_W S'
\]
qui, par définition de $W_S^u$ et $W_u^{S'}$, est compatible avec ces localiseurs
\[
\unsure{f^{*}|f}_{\underline{\mathrm{Fib}}_W}:(\underline{\mathrm{Fib}}_W S,W_S^u)\longrightarrow(\underline{\mathrm{Fib}}_W S',W_{S'}^u),
\]
et passe donc aux catégories de fractions en une flèche \unsure{(xx)|(**)}.

On peut se demander pourquoi ne pas être plus généreux, et définir $\mathrm{Ho}^{\unsure{\mathrm{lc}|\mathrm{le}}}_W(S)$ comme $(M_{/S})(W_S^u)^{-1}$, en définissant $W_S^u$ dans $M_{/S}$ de la \unsure{$=$|même} façon que tantôt dans $\underline{\mathrm{Fib}}_W S$ -- ce qui ne coûte pas plus cher. Une première réponse est d'expérience~: le localiseur $W_S^u$ dans $M_{/S}$ \deleted{tout} entier s'avère pratiquement inutilisable, d'une part parce que l'appartenance à $W_S^u$ (pour une flèche entre objets de $M_{/S}$ qui ne sont pas très particuliers, p. ex. $W$-fibrés sur $S$) est \deleted{\unsure{pratiquem}} invérifiable (sauf pour des iso et des choses comme ça ...), et que d'autre part la catégorie de fractions est prohibitivement
% archive sheet number: 261
\scanpage{262}
\mspage{212}

grosse et ne peut servir (semble-t-il) à rien. Une autre réponse, c'est qu' \deleted{\unsure{on}} \deleted{\unsure{doit}} pour la catégorie obtenue réponde bien à l'intuition de \deleted{\unsure{objets}} types d'homotopie relatifs \uline{localement constants} sur $S$ -- et ça veut dire que pour un chgt de base quelconque \inserted{$f:S'\to S$} par un $S'$ dans $M_{/S}$, d'où un foncteur composé
% slanted
\margin{$\unsure{\mathcal{F}|F}$ désigne ici une \unsure{ss/catégorie} fibrée de \deleted{\unsure{$\mathcal{F}$}} $M_{/M}=\underline{\mathrm{Fl}}(M)$ sur $M$, substitut éventuel de celle formée des $\underline{\mathrm{Fib}}_W\unsure{(S)|S}$, ou des $M_{/S}$ tout entiers}
\diagram{\begin{tikzcd}[ampersand replacement=\&,column sep=large]
\mathcal{F}_S(W_S^u)^{-1}\arrow[r,"{\deleted{\unsure{f^*}}\inserted{f_W^{*\,\unsure{\mathrm{lc}|\mathrm{le}}}}}"]\arrow[rr,rounded corners,to path={-- ([yshift=-1.25em]\tikztostart.south) -| (\tikztotarget.south)\tikztonodes},"{\unsure{\varphi|\phi}_{S'}}" near end,"{\deleted{\unsure{\mathcal{F}|F}}}"' very near end]
\&\mathcal{F}_{S'}(W_{S'}^u)^{-1}\arrow[r]\&\mathrm{Ho}_W
\end{tikzcd}}
déduit des foncteurs de catégories avec localiseurs
\diagram{\begin{tikzcd}[ampersand replacement=\&,column sep=large]
(\mathcal{F}_S,W_S^u)\arrow[r,"{\deleted{\unsure{\varphi_{S'}|i}} f^*}"]\arrow[rr,bend right=16,"{\deleted{\varphi_{S'}}}"']\arrow[rr,bend right=16,no head,decorate,decoration={zigzag,segment length=4pt,amplitude=1.6pt}]
\&(\mathcal{F}_{S'},W_{S'}^u)\arrow[r,"{\substack{\text{oubli}\\\text{de }S'}}"]\&(M,W)
\end{tikzcd}}
(foncteur \inserted{composé $\deleted{\unsure{\varphi_S}}\varphi_{S'}$} qui peut être vu comme le foncteur ``\uline{fibre homotopique en} \uline{l'objet $S'$ sur} $S$'', cet objet étant visualisé comme un ``point'' de $S$ ``à valeurs dans \unsure{$S'$|$S$}'' suivant le jargon fonctoriel général ...), le foncteur
\[
\begin{aligned}
S'&\longmapsto\varphi_{S'}(\unsure{\xi|\zeta})\qquad\text{pour }\xi\text{ dans }\deleted{\unsure{\mathrm{Ho}^{\mathrm{lc}}_W}}\mathcal{F}_S(W_S^u)^{-1}\\
M_{/S}&\longrightarrow\mathrm{Ho}_W
\end{aligned}
\]
transforme flèches de $W_S$ en isomorphismes~: la ``fibre $W$-homotopique'' en $S'/S$ d'un objet de $\mathcal{F}_S(W_S^u)^{-1}$ \deleted{doit être ``la} (vu comme \inserted{$W$-}type d'hom. \uline{loc. constant} sur $S$) doit être ``\unsure{la|le} \unsure{même}'', si $S'$ remplacé par $S''$ avec $S''\to S'$ donné et dans $W_S$. C'est pour cela qu'il est raisonnable de prendre $\mathcal{F}_S$ ``assez petit'' (savoir $\underline{\mathrm{Fib}}_W S$, qui est
% archive sheet number: 262
\scanpage{263}
\mspage{213}

exactement ajusté pour cela), plutôt que de prendre $M_{/S}$ tout entier.

On peut donc dire que la catégorie $\underline{\mathrm{Fib}}_W S$, et \unsure{son} localiseur $W_S^u$, ont été construits ad-hoc, le plus simplement possible, pour donner lieu à \unsure{un localisé|une localisée} $\mathrm{Ho}^{\unsure{\mathrm{lc}|\mathrm{le}}}_W(S)$ qui \deleted{\unsure{y}} réponde à l'intuition de \inserted{$W$-}types d'homotopie relatifs loc\textsuperscript{t} constants des deux façons suivantes~:

$1^{\circ})$ Il y a spontanément une loi contravariante en $S$, pour des flèches quelconques \deleted{\unsure{$f$:}} $S'\to S$ dans $\mathrm{Cat}$ (on a défini $W_S^u$ juste exprès en conséquence)

$2^{\circ})$ que la notion de ``\unsure{fibre|fibres} \inserted{$W$-}\unsure{homotopique|homotopiques} en $S'/S$» d'un objet $\unsure{\xi|\zeta}$ de $\mathrm{Ho}^{\mathrm{lc}}_W(S)$ satisfasse à la condition \unsure{précédente}, i.e. une flèche $S''\to S'$ dans $M_{/S}$ qui est dans $W_S$ définisse un \uline{iso} (dans $\mathrm{Ho}_W$) des fibres homotopiques
\diagram{\begin{tikzcd}[ampersand replacement=\&,row sep=small]
\xi_{S''}\arrow[r]\arrow[d,equal]\&\xi_{S'}\arrow[d,equal]\\
\unsure{\varphi|\phi}_{S''}(\xi)\&\varphi_{S'}(\xi)
\end{tikzcd}}
(et on a choisi $\unsure{\mathcal{F}|F}_{\unsure{C|S}}=\underline{\mathrm{Fib}}_W S$ juste exprès pour ça).

On a donc une excellente notion des fibres $W$-homotopiques pour des $X$ $W$-fibrés sur $S$ -- cette fois \inserted{excellente} pour des changements de base \uline{quelconques} -- et il n'y a \unsure{aucun mérite|aucune vérité} à ça, puisqu'on a tout défini exprès pour ça !
% archive sheet number: 263
\scanpage{264}
\mspage{214}

Pour le dire autrement, on a défini \uline{deux} notions de \inserted{$W$-}types d'homotopie relatifs sur $S$, l'une de façon à donner un tel type d'homotopie \inserted{$\mathrm{ho}^{!}_{W,S}(X)$} pour \uline{tout} objet de $M_{/S}$, mais il n'y a \inserted{a priori} une bonne notion de chgt de base que pour les chgt. de base $S'\unsure{\to|\text{---}}S$ qui sont $W$-fibrants; l'autre, qui correspond plus directement à l'intuition d'un $W$-type d'homotopie loc. constant sur $S$, n'associe un tel type $\mathrm{ho}^{\unsure{\mathrm{lc}|\mathrm{le}}}_{W,S}(X)$ qu'aux $X$ de $M_{/S}$ qui sont $W$-fibrés sur $S$, en revanche \unsure{cette} notion admet une bonne notion d'image inverse et de ``fibre $W$-homotopique'' pour des chgt. de base quelconques.

Une ``bonne'' théorie des fibres $W$-homotopiques doit réunir en une seule les deux notions précédentes, de sorte que les avantages des deux s'ajoutent, et les désavantages s'éliminent ! Elle s'obtient à l'aide d'un axiome draconien sur la situation (\unsure{alors|donc} que jusqu'à maintenant on n'a rien utilisé sur $W\subset\mathrm{Fl}(M)$, si ce n'est \deleted{\unsure{tout}} au plus la saturation faible) \unsure{,|;} en plus de l'existence des produits fibrés)~:
% archive sheet number: 264
\scanpage{265}
\mspage{215}

\uline{Axiome de la fibre homotopique} \textcircled{a} Pour tout objet $S$ dans $M$, le foncteur canonique
\[
\mathrm{Ho}^{\unsure{\mathrm{lc}|\mathrm{le}}}_W(S)\longrightarrow\mathrm{Ho}^{!}_W(S)
\]
déduit de l'inclusion
\[
(\underline{\mathrm{Fib}}_W S,W_S^u)\unsure{\text{---}|\longrightarrow|\longleftarrow}(M_{/S},W_S)
\]
est une \uline{équivalence de catégories}.

Cet axiome implique que dans $\underline{\mathrm{Fib}}_W S$, $W_S^u$ et le localiseur induit par $W_S$ (que je désigne encore par $W_S$ par abus de notation) ont \unsure{$=$|même} saturé fort. En fait, \deleted{\unsure{l'axio}} il est naturel de compléter l'axiome précédent par le suivant

\textcircled{b} (Axiome adjoint de l'axiome de la fibre homotopique) dans $\underline{\mathrm{Fib}}_W S$, les localiseurs $W_S^u$ et $W_S$ coïncident, i.e. on a ceci~:

Soient $X,Y$ dans $\underline{\mathrm{Fib}}_W S$, $X\xrightarrow{u}Y$ un $S$-morphisme tel que $u\in W_S$, alors $u\in W_S^u$ (i.e. le morphisme \inserted{$u':X'\unsure{\text{---}|\to}Y'$} déduit par \unsure{tt} changement de base $S'\to S$ est encore dans $W$).
% archive sheet number: 265
\scanpage{266}
\mspage{216}

Je présume que a) sera rarement satisfait sans que \unsure{b|S} le soit (et inversement d'ailleurs), et que la démonstration de a) et b) se fera pratiquement en une seule haleine.

Si $M=\mathrm{Cat}$, $W$ un localiseur fondamental, alors la \deleted{\unsure{validité de a)}} notion adoptée ici de $W$-fibration (lire~: ``ultra forte'', dans la terminologie spécifique développée dans $\mathrm{Cat}$) est à tel point forte, qu'elle implique déjà \inserted{pratiquement} b) -- (moyennant $W(4\unsure{a'|a})$, il est vrai, cf p. 53), et a) n'est guère vérifiable que si on a $W(4a)^{*}$ --
% slanted
\margin{Non, il me semble que $W(4b)$ doit suffire ...}
(loc. cit.). Donc ici c'est ``l'axiome des fibrations'' sur $W$, sous sa forme forte $W(4a)$, qui est la clef de voûte d'une bonne théorie des fibres homotopiques.

Une telle théorie (revenant au cas général) revient à dire, \inserted{pour l'essentiel,} qu'il y a ``assez'' de $W$-fibrations (``ultra fortes'') pour pouvoir décrire \uline{tout} objet d'un $M_{/S}$, à $W_S$-équivalence près, comme un objet $W$-fibrant sur $S$. Dans \deleted{\unsure{tous}} les cas à ma connaissance, on a \unsure{de|bien|alors} plus ceci~: on
% archive sheet number: 266
